\section{Objetivos}

Que el alumno:
\begin{itemize}
	\item Analizar en régimen senoidal permanente un sistema trifásico alimentado por un STB.
	\item Verificar mediante fasores el desfase aproximado de \SI{120}{\degree} entre las tres fases.
	\item Comparar voltajes de fase y de línea, validando la relación $V_L=\sqrt{3}V_F$.
	\item Analizar corrientes de fase y de línea en cargas conectadas en estrella y delta.
	\item Estudiar el efecto del desbalance por pérdida de fase, impedancias distintas y neutro flotante.
	\item Complementar el análisis con simulaciones, oscilogramas y diagramas fasoriales elaborados en Falstad.
\end{itemize}

\section{Equipo y material empleado}
\begin{table}[H]
	\centering
	\caption{Material empleado durante la práctica.}
	\begin{tabular}{clc}
		\toprule
		Cantidad & Elemento & Valor nominal \\
		\midrule
		3 & Resistencias & \SI{22}{\kilo\ohm}, \SI{0.5}{W} \\
		3 & Resistencias & \SI{1}{\kilo\ohm}, \SI{0.5}{W} \\
		3 & Resistencias & \SI{33}{\ohm}, \SI{0.5}{W} \\
		1 & Capacitor & \SI{0.22}{\micro F} \\
		1 & Simulador trifásico balanceado & STB \\
		1 & Osciloscopio / simulador Falstad & Medición de formas de onda \\
		\bottomrule
	\end{tabular}
\end{table}